\section*{الفصل \ref{ch:b2:surfaces} --- السطوح}

\begin{solution}{exo:b2:surfaces:1}
مع $\sigma(\theta, r) = (r\cos\theta,\ r\sin\theta,\ r)$:
\[
\sigma_\theta = (-r\sin\theta,\ r\cos\theta,\ 0),
\qquad
\sigma_r = (\cos\theta,\ \sin\theta,\ 1),
\]
وهما مستقلان من أجل $r > 0$. \omterm{def:b2:surfaces:tangent}{والمستوي المماس} في $M_0 =
\sigma(\theta_0, r_0)$ يمر بالنقطة $M_0$
باتجاهي $\sigma_\theta, \sigma_r$. والآن فإن $M_0 - O = r_0(\cos\theta_0,
\sin\theta_0, 1) = r_0\,\sigma_r(\theta_0, r_0)$ نفسه اتجاه \omterm{def:b2:curves:arc}{مماس}: ومنه فإن الرأس $O$
يقع على \omterm{def:b2:surfaces:tangent}{المستوي المماس}. (وهذا هو السلوك العام للمخاريط: فهي مسطّرة
بمستقيمات تمر بالرأس، ويحوي \omterm{def:b2:surfaces:tangent}{المستوي المماس} المولّد المستقيمي المار
بنقطة التماس.)
\end{solution}

\begin{solution}{exo:b2:surfaces:2}
نطبق \cref{prop:b2:surfaces:gradient} على $f(x,y,z) =
\frac{x^2}{a^2} + \frac{y^2}{b^2} + \frac{z^2}{c^2}$: $\nabla
f(x_0,y_0,z_0) = 2\bigl(\frac{x_0}{a^2}, \frac{y_0}{b^2},
\frac{z_0}{c^2}\bigr) \neq 0$ على السطح. \omterm{def:b2:surfaces:tangent}{والمستوي المماس} هو
\[
\frac{x_0}{a^2}(x - x_0) + \frac{y_0}{b^2}(y - y_0)
+ \frac{z_0}{c^2}(z - z_0) = 0,
\qquad\text{أي}\qquad
\frac{x_0\,x}{a^2} + \frac{y_0\,y}{b^2} + \frac{z_0\,z}{c^2} = 1,
\]
باستعمال كون $(x_0, y_0, z_0)$ يحقق معادلة المجسم الإهليلجي --- وهي قاعدة ``شطر
المربعات'' التي تعمّم المقدار $\langle
M_0, M \rangle = R^2$ في حالة الكرة.
\end{solution}

\begin{solution}{exo:b2:surfaces:3}
$\sigma_u = (-v\sin u,\ v\cos u,\ a)$ و$\sigma_v = (\cos u,\
\sin u,\ 0)$، ومنه
\[
E = v^2 + a^2, \qquad F = 0, \qquad G = 1,
\qquad
\sqrt{EG - F^2} = \sqrt{v^2 + a^2} > 0
\]
(فالسطح اللولبي نظامي في كل نقطة، بما في ذلك على محوره $v =
0$).
\omterm{def:b2:surfaces:area}{ومساحة} القطعة:
\[
\mathcal{A} = \int_0^{2\pi}\!\!\int_0^1
\sqrt{v^2 + a^2}\;\dd v\,\dd u
= 2\pi\cdot\frac12\Bigl[v\sqrt{v^2 + a^2}
+ a^2\ln\bigl(v + \sqrt{v^2 + a^2}\bigr)\Bigr]_0^1 ,
\]
أي $\mathcal{A} = \pi\Bigl(\sqrt{1 + a^2} +
a^2\ln\frac{1 + \sqrt{1 + a^2}}{a}\Bigr)$.
\end{solution}

\begin{solution}{exo:b2:surfaces:4}
المشتقات:
\[
\sigma_\theta = \bigl(-(R + r\cos\psi)\sin\theta,\
(R + r\cos\psi)\cos\theta,\ 0\bigr),
\qquad
\sigma_\psi = \bigl(-r\sin\psi\cos\theta,\
-r\sin\psi\sin\theta,\ r\cos\psi\bigr).
\]
ثم
\[
E = (R + r\cos\psi)^2,
\qquad
F = r\sin\psi\cos\psi\,(R + r\cos\psi)
\bigl(\sin\theta\cos\theta - \sin\theta\cos\theta\bigr) = 0,
\qquad
G = r^2 ,
\]
ومنه $\sqrt{EG - F^2} = r(R + r\cos\psi) \geq r(R - r) > 0$: فهو نظامي في كل نقطة. \omterm{def:b2:surfaces:area}{والمساحة}:
\[
\mathcal{A}
= \int_0^{2\pi}\!\!\int_0^{2\pi} r(R + r\cos\psi)
\,\dd\theta\,\dd\psi
= 2\pi r \Bigl(2\pi R + r\int_0^{2\pi}\cos\psi\,\dd\psi\Bigr)
= 4\pi^2 R r ,
\]
مع تكامل الحد $\cos\psi$ إلى صفر --- وهي مبرهنة بابوس: \omterm{def:b2:surfaces:area}{المساحة} تساوي
$=$ (طول الدائرة المدارة) في $\times$ (المسافة التي يقطعها مركزها).
\end{solution}

\begin{solution}{exo:b2:surfaces:5}
من أجل $\sigma(x, y) = (x, y, f(x,y))$: $\sigma_x = (1, 0, f_x)$، $\sigma_y = (0, 1, f_y)$، ومنه $E = 1 + f_x^2$، $F = f_xf_y$، $G = 1
+ f_y^2$ و
\[
EG - F^2 = (1 + f_x^2)(1 + f_y^2) - f_x^2f_y^2
= 1 + f_x^2 + f_y^2 ,
\]
وهو ما يعطي صيغة \omterm{def:b2:surfaces:area}{المساحة} المذكورة. ومن أجل $f = \frac12(x^2 + y^2)$ على قرص الوحدة:
$1 + f_x^2 + f_y^2 = 1 + x^2 + y^2$، وبالإحداثيات القطبية ($x = \rho\cos\alpha$، $y = \rho\sin\alpha$، والجاكوبي $\rho$،
\cref{ch:b2:multint}):
\[
\mathcal{A}
= \int_0^{2\pi}\!\!\int_0^1 \sqrt{1 + \rho^2}\;\rho
\,\dd\rho\,\dd\alpha
= 2\pi\Bigl[\tfrac13(1 + \rho^2)^{3/2}\Bigr]_0^1
= \frac{2\pi}{3}\bigl(2\sqrt2 - 1\bigr).
\]
\end{solution}

\begin{solution}{exo:b2:surfaces:6}
انطلاقا من حساب \cref{ex:b2:surfaces:spherearea}، $E =
R^2\cos^2\varphi$، $F = 0$، $G = R^2$، ومنه بحسب \cref{prop:b2:surfaces:curvelength}
\[
L = \int_a^b R\sqrt{\cos^2\varphi(t)\,\theta'(t)^2
+ \varphi'(t)^2}\;\dd t .
\]
وليكن الطرفان $(\theta_0, \varphi_1)$ و$(\theta_0,
\varphi_2)$، $\varphi_1 < \varphi_2$. ومن أجل أي منحن واصل بينهما،
\[
L \geq \int_a^b R\,\abs{\varphi'(t)}\,\dd t
\geq R\,\Bigl|\int_a^b \varphi'(t)\,\dd t\Bigr|
= R\,(\varphi_2 - \varphi_1),
\]
بإهمال الحد الموجب $\cos^2\varphi\,\theta'^2$ واستعمال متراجحة المثلث في التكاملات. أما
قوس خط الزوال $\theta \equiv \theta_0$، مع $\varphi$ متزايدة من $\varphi_1$ إلى $\varphi_2$، فطوله
$R(\varphi_2 - \varphi_1)$ بالضبط: فهو الأقصر. (وخطوط الزوال دوائر كبرى؛ وهذه هي الحالة
الأولى الأولية من كون جيوديزيات الكرة دوائر كبرى.)
\end{solution}

\begin{solution}{exo:b2:surfaces:7}
نثبّت منحنى $\gamma$ مرسوما على $S$ ونضع $g(t) = \norm{\gamma(t)
- \Omega}^2$. عندئذ $g'(t) = 2\langle \gamma'(t),\ \gamma(t) -
\Omega\rangle$. ويمر
الناظم في $M = \gamma(t)$ بالنقطة $\Omega$ بحكم الفرضية، ومنه فإن $\gamma(t) - \Omega$ متجهة
\emph{ناظمية}، متعامدة مع \omterm{def:b2:surfaces:tangent}{المستوي المماس}، وبوجه خاص مع السرعة $\gamma'(t)$
(\cref{prop:b2:surfaces:velocity}): ومنه $g' =
0$، ويكون $g$ ثابتا على طول كل منحن مرسوم.

والآن فإن المجموعة $S_c = \{M \in S : \norm{M - \Omega}^2 = c\}$ مغلقة في $S$؛ وهي أيضا \omterm{def:b2:metric:topology}{مفتوحة} في $S$:
فبجوار أي نقطة منها يكون $S$ تمثيلا بيانيا نظاميا، ومنه فإن أي
نقطة مجاورة من $S$ تُوصل بها بمنحن مرسوم (قطعة مرفوعة)، يكون
$g$ ثابتا على طوله. وبما أن $S$ مترابط و$S_c$ غير خالية من
أجل $c$ المناسبة، فإن $S = S_c
\subseteq$ هي الكرة ذات المركز $\Omega$ ونصف
القطر $\sqrt c$ (\cref{ch:b2:metric}: بحجة الترابط).
\end{solution}

\begin{solution}{exo:b2:surfaces:8}
نكتب $\Phi(u', v') = (u, v)$. وبقاعدة السلسلة،
\[
\tilde\sigma_{u'} = \frac{\partial u}{\partial u'}\sigma_u
+ \frac{\partial v}{\partial u'}\sigma_v,
\qquad
\tilde\sigma_{v'} = \frac{\partial u}{\partial v'}\sigma_u
+ \frac{\partial v}{\partial v'}\sigma_v ,
\]
حيث تُقوَّم المشتقات الجزئية للتطبيق $\sigma$ في $\Phi(u',v')$. وبنشر الجداء
الشعاعي بثنائية الخطية واستعمال $\sigma_u \wedge \sigma_u =
\sigma_v \wedge \sigma_v = 0$ و$\sigma_v \wedge \sigma_u =
-\sigma_u \wedge \sigma_v$:
\[
\tilde\sigma_{u'} \wedge \tilde\sigma_{v'}
= \Bigl(\frac{\partial u}{\partial u'}
\frac{\partial v}{\partial v'}
- \frac{\partial v}{\partial u'}
\frac{\partial u}{\partial v'}\Bigr)\,
\sigma_u \wedge \sigma_v
= \det J_\Phi \cdot (\sigma_u \wedge \sigma_v)\circ\Phi .
\]
وبأخذ المعايير نحصل على المتطابقة. ثم، بصيغة تغيير المتغيرات
(\cref{ch:b2:multint}) مطبقة على التطبيق $\Phi$ على $K' =
\Phi^{-1}(K)$:
\[
\iint_{K'} \norm{\tilde\sigma_{u'} \wedge \tilde\sigma_{v'}}
\,\dd u'\dd v'
= \iint_{K'} \norm{\sigma_u \wedge \sigma_v}\circ\Phi\;
\abs{\det J_\Phi}\,\dd u'\dd v'
= \iint_K \norm{\sigma_u \wedge \sigma_v}\,\dd u\,\dd v :
\]
فالتوسيمان يعطيان \omterm{def:b2:surfaces:area}{المساحة} نفسها للقطعة نفسها من السطح.
\end{solution}

\begin{solution}{exo:b2:surfaces:9}
في الخريطة الكروية، $z = R\sin\varphi$، ويوافق النطاق $\varphi_1 \leq \varphi \leq \varphi_2$ مع $z_i = R\sin\varphi_i$. ومع
عنصر \omterm{def:b2:surfaces:area}{المساحة}
$R^2\cos\varphi\,\dd\theta\,\dd\varphi$
(\cref{ex:b2:surfaces:spherearea}):
\[
\mathcal A = \int_{\varphi_1}^{\varphi_2}\!\!\int_0^{2\pi}
R^2\cos\varphi\,\dd\theta\,\dd\varphi
= 2\pi R^2(\sin\varphi_2 - \sin\varphi_1)
= 2\pi R\,(z_2 - z_1) .
\]
ولا تتعلق النتيجة إلا بالارتفاع $z_2 - z_1$: فللشرائح المتساوية السماكة
مساحات متساوية، سواء قُطعت عند خط الاستواء أو عند القطب --- وهي
مبرهنة أرخميدس في علبة القبعة، والسبب في أن \omterm{def:b2:surfaces:area}{المساحة} الجانبية
للأسطوانة المحيطة ($2\pi R \cdot
2R = 4\pi R^2$) تساوي \omterm{def:b2:surfaces:area}{مساحة} الكرة.
\end{solution}

\begin{solution}{exo:b2:surfaces:10}
$\sigma_\theta = (-r\sin\theta,\ r\cos\theta,\ 0)$ و$\sigma_z = (r'\cos\theta,\ r'\sin\theta,\ 1)$، ومنه
\[
\sigma_\theta \wedge \sigma_z
= (r\cos\theta,\ r\sin\theta,\ -r\,r') ,
\]
وهي متجهة ناظمية في $M = (r\cos\theta, r\sin\theta, z)$. والناظم هو
\[
t \mapsto \bigl(r(1 + t)\cos\theta,\ r(1 + t)\sin\theta,\
z - t\,r\,r'\bigr),
\]
وهو يبلغ عند $t = -1$ النقطة $(0,\ 0,\ z + r(z)\,r'(z))$: فكل ناظم يلاقي المحور، عند
الارتفاع $z + rr'$. (وهذا هو السبب ذو الأبعاد الثلاثة في بقاء التناظر
الدوراني في حقل النواظم.)
\end{solution}

\begin{solution}{exo:b2:surfaces:11}
$\sigma_u = (-\sin u, \cos u, 0)$، $\sigma_v = (0, 0, 1)$: $E
= 1$، $F = 0$، $G = 1$. وبحسب \cref{prop:b2:surfaces:curvelength}، يكون طول
منحن مرسوم $\int\sqrt{u'^2 + v'^2}\,\dd t$ --- أي طول ظله الوسيطي $(u(t), v(t))$ في المستوي: فالخريطة
تقايس محلي (وهو فرد الأسطوانة). واللولب $t \mapsto (\cos t, \sin t, ct)$، $t \in \intcc0{2\pi}$، ظله القطعة
الواصلة من $(0,0)$ إلى $(2\pi, 2\pi c)$، وطولها $2\pi\sqrt{1 + c^2}$. وأي منحن مرسوم من
$(1,0,0)$ إلى $(1,
0, 2\pi c)$ يصنع دورة كاملة واحدة له ظل \omterm{def:b2:metric:continuity}{متصل} يصل $(0, 0)$
بالنقطة $(2\pi, 2\pi c)$، وطوله المستوي $\geq$ القطعة المستقيمة؛ وبما أن
الطولين متساويان، فإن اللولب هو الأقصر.
\end{solution}

\begin{solution}{exo:b2:surfaces:12}
الدالة $g(M) = \norm{M}^2$ متصلة على المتراص $S$، ومنه فهي تبلغ قيمتها
العظمى في نقطة $M_0$. ومن أجل كل منحن $\gamma$ مرسوم على $S$ مع
$\gamma(0) = M_0$، تبلغ الدالة $t
\mapsto \norm{\gamma(t)}^2$ قيمة عظمى في $t = 0$، ومنه تنعدم مشتقتها
$2\langle\gamma'(0), M_0\rangle$: أي إن $M_0$ متعامد مع كل متجهة مماسة، أي إنه ناظمي للسطح
$S$ في $M_0$. وبما أن $\nabla f(M_0) \neq 0$ يوجه الناظم أيضا (\cref{prop:b2:surfaces:gradient})، فإن
$\nabla f(M_0)$ و$\vect{OM_0}$ خطيان. وفي حالة المجسم الإهليلجي، تعني القطرية
\[
\Bigl(\frac{x_0}{a^2}, \frac{y_0}{b^2},
\frac{z_0}{c^2}\Bigr) = \mu\,(x_0, y_0, z_0) :
\]
فكل إحداثية تحقق $x_0(\frac1{a^2} - \mu) = 0$، وهكذا؛ وبما أن المقادير $a^{-2}, b^{-2}, c^{-2}$ متمايزة،
فإن إحداثية واحدة على الأكثر غير معدومة، والحلول على السطح هي أطراف
المحاور الستة $(\pm a, 0, 0)$ و$(0, \pm b, 0)$ و$(0,
0, \pm c)$. والنقطتان الأبعد هما $(\pm a, 0, 0)$،
على المسافة $a = \max(a,b,c)$.
\end{solution}

\begin{solution}{pb:b2:surfaces:1}
\textbf{1.} بالنشر، وباستعمال تناظر $A$ ($t^{\mathsf T}\!APY = (At)^{\mathsf T}PY$):
\[
q(PY + t) = Y^{\mathsf T}P^{\mathsf T}\!APY +
2\,(At + b)^{\mathsf T}PY + \bigl(t^{\mathsf T}\!At +
2b^{\mathsf T}t + c\bigr),
\]
وهي الثلاثية المكتوبة. والمصفوفة $A' = P^{\mathsf T}\!AP =
P^{-1}AP$ مشابهة للمصفوفة $A$: فلهما
\omterm{def:b2:reduction:charpoly}{كثير الحدود المميز} نفسه، \omterm{def:b2:reduction:eigen}{والطيف} والرتبة والبصمة أنفسها. وأخيرا
$\{q = 0\} = \{sq = 0\}$ من أجل $s \neq 0$، ومنه فإن ما يرتبط بالمجموعة هو المعادلة
\emph{إلى غاية عدد} فحسب؛ والضرب في $s$ يضرب جميع \omterm{def:b2:reduction:eigen}{القيم الذاتية}
في $s$.

\textbf{2.} توفر المبرهنة الطيفية مصفوفة $P \in O(3)$ تحقق $P^{\mathsf T}\!AP = \operatorname{diag}(\lambda_1, \lambda_2,
\lambda_3)$؛
ويحول السؤال 1 مع $t = 0$ المعادلة إلى $\sum\lambda_iy_i^2 + 2\sum\beta_iy_i + c = 0$، حيث $\beta =
P^{\mathsf T}b$.

\textbf{3.} من أجل $\lambda_i \neq 0$: $\lambda_iy_i^2 +
2\beta_iy_i = \lambda_i\bigl(y_i +
\frac{\beta_i}{\lambda_i}\bigr)^2 -
\frac{\beta_i^2}{\lambda_i}$؛ ويعطي الانسحاب $z_i = y_i +
\beta_i/\lambda_i$ (و$z_i = y_i$
من أجل $i > r$)
\[
\sum_{i \leq r}\lambda_iz_i^2 + 2\sum_{i > r}\beta_iz_i +
c'' = 0, \qquad
c'' = c - \sum_{i\leq r}\frac{\beta_i^2}{\lambda_i}.
\]

\textbf{4.} $q(2\Omega - X) = (2\Omega - X)^{\mathsf T}
A(2\Omega - X) + 2b^{\mathsf T}(2\Omega - X) + c$؛ وبالنشر وطرح $q(X)$، تتلاشى الحدود التربيعية
ويبقى
\[
q(2\Omega - X) - q(X) = -4\,\langle A\Omega + b,\ X\rangle +
4\,\langle A\Omega + b,\ \Omega\rangle .
\]
فإذا كان $A\Omega + b = 0$ انعدم هذا تطابقا: ومنه فإن التناظر المركزي يحافظ على
$q$، ومنه على \omterm{pb:b2:surfaces:1}{السطح من الدرجة الثانية}. وبالعكس، فإن القول
``$q(2\Omega - X) = q(X)$ من أجل كل $X$'' يعني أن الدالة الأفينية أعلاه تنعدم على
$\R^3$ كله، وهو ما يفرض انعدام جزئها الخطي $A\Omega + b$. ومنه فإن المراكز
$=$ هي حلول $A\Omega = -b$: وهي مجموعة غير خالية إذا وفقط إذا كان $b \in \operatorname{im}A$
(وهي \omterm{def:b2:affine:def}{فضاء أفيني} جزئي موجه بالفضاء $\ker A$)، ونقطة وحيدة إذا وفقط إذا
كانت $A$ قابلة للقلب.

\textbf{5.} البصمة $(3,0)$ (وجميع المقادير $\lambda_i > 0$): يعطي $\delta >
0$
المقدار $\frac{x^2}{a^2} + \frac{y^2}{b^2} + \frac{z^2}{c^2}
= 1$ مع $a = \sqrt{\delta/\lambda_1}$، وهكذا --- أي مجسم إهليلجي؛ و$\delta = 0$:
النقطة الوحيدة $O$؛ و$\delta < 0$: خال. والبصمة $(2,1)$ ($\lambda_1, \lambda_2 > 0 >
\lambda_3$):
$\delta > 0$: $\frac{x^2}{a^2} + \frac{y^2}{b^2}
- \frac{z^2}{c^2} = 1$، وهو المجسم الزائدي ذو الطية الواحدة؛ و$\delta =
0$:
المخروط $\frac{x^2}{a^2} + \frac{y^2}{b^2} =
\frac{z^2}{c^2}$؛ و$\delta < 0$: $\frac{z^2}{c^2} -
\frac{x^2}{a^2} - \frac{y^2}{b^2} = 1$، وهو المجسم الزائدي ذو الطيتين
($\abs z \geq c$: بمركبتين).

\textbf{6.} للجزء التربيعي المصفوفة $A$ التي قطرها $1$ وباقي
معاملاتها $2$: $A = 2J - I$. وبما أن \omterm{def:b2:reduction:eigen}{طيف} $J$ هو $\{3, 0, 0\}$ (\omterm{def:b2:reduction:eigen}{بالمتجهة
الذاتية} $(1,1,1)$ من أجل $3$)، فإن \omterm{def:b2:reduction:eigen}{طيف} $A$ هو $\{5, -1, -1\}$، \omterm{def:b2:reduction:eigen}{والقيمة
الذاتية} $5$ تحملها المتجهة $\R(1,1,1)$. وفي الإحداثيات المدارة: $5u^2 - v^2 - w^2 =
1$،
أي $\frac{u^2}{1/5} - v^2 - w^2 = 1$: وهو مجسم زائدي ذو طيتين، ودوراني (فالقيمتان الذاتيتان
$-1$ متساويتان) حول المحور $\R(1,1,1)$.

\textbf{7.} السطح هو $\bigl(\frac xa - \frac
zc\bigr)\bigl(\frac xa + \frac zc\bigr) = \bigl(1 - \frac
yb\bigr)\bigl(1 + \frac yb\bigr)$. ومن أجل $(\lambda : \mu) \neq
(0:0)$ نعرّف المستقيم
\[
D_{\lambda:\mu} :\quad
\lambda\Bigl(\frac xa - \frac zc\Bigr) = \mu\Bigl(1 - \frac
yb\Bigr), \qquad
\mu\Bigl(\frac xa + \frac zc\Bigr) = \lambda\Bigl(1 + \frac
yb\Bigr)
\]
(بمعادلتين أفينيتين مستقلتين: أي مستقيم). ويبين ضرب المعادلتين أن
كل نقطة من $D_{\lambda:\mu}$ تقع على السطح عندما $\lambda\mu \neq 0$؛ أما الحالتان $\lambda = 0$
و$\mu = 0$ فيُتحقق منهما مباشرة (مثلا $\lambda = 0$: $y =
b$، $\frac xa = -\frac zc$، وهو
يحقق المعادلة). والعائلة الثانية $D'_{\lambda:\mu}$ تبادل العاملين في الطرف
الأيمن.

\textbf{8.} نثبّت $M$ على السطح. وتشكل شروط $M \in
D_{\lambda:\mu}$ جملة خطية
متجانسة من الحجم $2\times2$ بدلالة $(\lambda, \mu)$ محددها
\[
\Bigl(\frac{x^2}{a^2} - \frac{z^2}{c^2}\Bigr) - \Bigl(1 -
\frac{y^2}{b^2}\Bigr) = 0
\]
وذلك بالضبط لأن $M$ يقع على \omterm{pb:b2:surfaces:1}{السطح من الدرجة الثانية}: ومنه يوجد
حل غير تافه $(\lambda : \mu)$. ومصفوفة المعاملات ليست معدومة أبدا (وإلا فرض
ذلك $1 - \frac yb = 1 + \frac yb =
0$)، ومنه فرتبتها $1$ والحل وحيد إلى غاية عدد: أي إن
مستقيما واحدا بالضبط من العائلة يمر بالنقطة $M$. والأمر نفسه
يصح للعائلة الثانية، والمستقيمان متمايزان (فعند $(a, 0, 0)$ هما $\{x = a,\ \frac zc = \frac
yb\}$
و$\{x = a,\ \frac zc = -\frac yb\}$): ومنه فإن المجسم الزائدي ذا الطية الواحدة مزدوج التسطير.

\textbf{9.} لتكن $M = (x, y, z)$ على المجسم الزائدي، $\rho^2 =
\frac{x^2}{a^2} + \frac{y^2}{b^2} = 1 + \frac{z^2}{c^2} \geq
1$. والنقطة $N = (x,\ y,\ \varepsilon c\rho)$
مع $\varepsilon$ إشارة $z$ تحقق $\frac{x^2}{a^2} +
\frac{y^2}{b^2} - \frac{(c\rho)^2}{c^2} = 0$: $N \in C$، و
\[
d(M, C) \leq \abs{z - \varepsilon c\rho}
= c\bigl(\rho - \sqrt{\rho^2 - 1}\bigr)
= \frac{c}{\rho + \sqrt{\rho^2 - 1}} .
\]
فإذا كان $\norm M \to \infty$ فإن $\rho \to \infty$ (فالإحداثيات الثلاث محصورة بمضاعفات
$\rho$)، ومنه $d(M, C)
\to 0$. والمقطع $x = a$: $\frac{y^2}{b^2} -
\frac{z^2}{c^2} = 0$، وهو زوج المستقيمين
المتقاطعين في السؤال 8 --- وكذلك عند $x = -a$.

\textbf{10.} مع $\lambda_3 = 0$، يترك السؤال 3 المعادلة $\lambda_1z_1^2 + \lambda_2z_2^2 + 2\beta_3z_3 + c'' = 0$. وفي الأساس
الذاتي، $\operatorname{im}A = \operatorname{Vect}(e_1,
e_2)$، ومنه بحسب السؤال 4 توجد مراكز إذا وفقط إذا كان
$\beta_3 = 0$. فإذا كان $\beta_3 \neq 0$: فإن الانسحاب $z_3 \mapsto z_3 -
c''/(2\beta_3)$ يزيل الثابت، ويبقى
$\lambda_1z_1^2
+ \lambda_2z_2^2 + 2\beta_3z_3 = 0$، أي $z_3 = px^2 +
qy^2$ بعد إعادة التسمية: أي \emph{مجسم مكافئ إهليلجي}
إذا كان $\lambda_1\lambda_2 > 0$، و\emph{مجسم مكافئ زائدي} إذا كان $\lambda_1\lambda_2 < 0$ --- ولا
مركز فعلا. وإذا كان $\beta_3 = 0$: فالمعادلة $\lambda_1z_1^2 + \lambda_2z_2^2 +
c'' = 0$ لا تتضمن $z_3$: ومنه
فإن السطح أسطوانة فوق المقطع المخروطي المستوي الموافق --- أسطوانة
إهليلجية أو مستقيم أو مجموعة خالية عندما $\lambda_1\lambda_2 > 0$؛ وأسطوانة زائدية
أو زوج مستويين متقاطعين عندما $\lambda_1\lambda_2 < 0$ --- مع مستقيم كامل من المراكز
$\{(z_1^*, z_2^*)\} \times \R$.

\textbf{11.} $xy - z = 0$. وبالتعويض $x = \frac{u +
v}{\sqrt2}$، $y = \frac{u - v}{\sqrt2}$ (وهو دوران بالزاوية
$\pi/4$): $xy = \frac{u^2 - v^2}2$، ومنه تصير المعادلة $z =
\frac12(u^2 - v^2)$: أي مجسم مكافئ زائدي
--- وهو سرج الشكل، إلى غاية العامل $\frac12$.

\textbf{12.} المستقيم $\{x = x_0,\ z = x_0y\}$ (الموسّم بالوسيط $y$) يقع بوضوح على
$z = xy$، وكذلك $\{y = y_0,\ z =
xy_0\}$؛ ويمر بالنقطة $(x_0, y_0, x_0y_0)$ كلاهما.
والوحدانية: إذا بقي $t \mapsto (x_0 + tv_1, y_0 + tv_2, z_0 +
tv_3)$ على السطح، أعطى معامل $t^2$ في $(x_0
+ tv_1)(y_0 + tv_2) - z_0 - tv_3$
المقدار $v_1v_2 = 0$، ومنه $v_1 = 0$ أو $v_2 = 0$، فنحط في إحدى العائلتين:
أي مستقيم واحد من كل عائلة بكل نقطة --- وهو ثاني \omterm{pb:b2:surfaces:1}{سطح من الدرجة
الثانية} مزدوج التسطير.

\textbf{13.} بإتمام المربعات: $(x - 1)^2 + (y + 2)^2 = 2$، دون أي شرط على $z$: أي
أسطوانة دائرية قائمة نصف قطرها $\sqrt2$ ومحورها المستقيم الشاقولي
$\{(1, -2, z) : z \in
\R\}$ --- أي مستقيم من المراكز، كما يتنبأ السؤال 10.

\textbf{14.} مع $\lambda_2 = \lambda_3 = 0$ تكون المعادلة المختزلة $\lambda_1z_1^2 + 2\beta_2z_2 + 2\beta_3z_3 + c''
= 0$. ودوران لمستوي
النواة $(z_2, z_3)$ يوائم الصورة الخطية: $2\beta_2z_2 + 2\beta_3z_3 = 2\beta w$ مع $\beta = \sqrt{\beta_2^2 + \beta_3^2}$. فإذا كان
$\beta \neq 0$، أزح $w$ لامتصاص $c''$: $\lambda_1z_1^2 + 2\beta w =
0$، وهي \emph{أسطوانة
مكافئة}؛ وإذا كان $\beta = 0$: فإن $\lambda_1z_1^2 = -c''$ يعطي مستويين متوازيين ($c''
\lambda_1 < 0$)،
أو مستويا مضاعفا ($c'' = 0$)، أو المجموعة الخالية. ومن أجل $(x + y)^2 = z$: مع
$u = \frac{x + y}{\sqrt2}$ تُكتب المعادلة $z = 2u^2$: أي أسطوانة مكافئة، ثابتة بالانسحابات
على طول $(1, -1, 0)$.

\textbf{15.} تُرسل كل \omterm{pb:b2:surfaces:1}{سطح من الدرجة الثانية} بدوران زائد انسحابات على
أحد ما يلي: (الرتبة 3) مجسم إهليلجي، نقطة، مجموعة خالية، مجسم زائدي
ذو طية واحدة، مخروط، مجسم زائدي ذو طيتين؛ (الرتبة 2) مجسم مكافئ
إهليلجي، مجسم مكافئ زائدي، أسطوانة إهليلجية، مستقيم، مجموعة خالية،
أسطوانة زائدية، زوج مستويين متقاطعين؛ (الرتبة 1) أسطوانة مكافئة، زوج
مستويين متوازيين، مستو مضاعف، مجموعة خالية. وبعدّ المتغيرات الخالية
الثلاثة \emph{أنماطا أفينية} متمايزة من المعادلات، يكون العدد سبعة
عشر؛ ومن بينها تسعة سطوح حقيقية: المجسم الإهليلجي، والمجسمان
الزائديان، والمخروط، والمجسمان المكافئان، والأسطوانات الثلاث.

\textbf{16.} الخوارزمية. (1) اقرأ $(A, b, c)$؛ واحسب \omterm{def:b2:reduction:charpoly}{كثير الحدود المميز}
للمصفوفة $A$ وقيمها الذاتية (وهي حقيقية، بالمبرهنة الطيفية)
والمقدار $r = \operatorname{rank}A$. (2) حل $A\Omega = -b$ (بحذف غاوس): فهي قابلة للحل أو غير
قابلة --- أي مراكز أو لا مراكز. (3) وإذا كانت قابلة للحل، فانسحب إلى
مركز: تصير المعادلة $\sum\lambda_iz_i^2 + c'' = 0$ مع $c''
= c + \langle b, \Omega\rangle$؛ ثم رتّب بحسب $r$ والبصمة
وإشارة $c''$ باستعمال الأسئلة 5 و10 و14. (4) وإذا لم تكن قابلة
للحل ($r \leq 2$)، فأدر واختزل كما في السؤالين 10 و14: مجسم مكافئ
($r = 2$) أو أسطوانة مكافئة ($r = 1$)، إهليلجي أو زائدي بحسب إشارة
$\lambda_1\lambda_2$. وكل خطوة حساب منته: جذور معادلة تكعيبية جذورها حقيقية،
ورتب، وجمل خطية.

\textbf{17.} المصفوفة $A$ قطرها $1$ وباقي معاملاتها $-1$:
$A = 2I
- J$، وطيفها $\{2 - 3,\ 2,\ 2\} = \{-1, 2, 2\}$ مع $-1$ على $\R(1,1,1)$. والبصمة $(2,1)$،
$b = 0$، والطرف الأيمن $\delta = 1 > 0$: $2u^2 + 2v^2 - w^2 = 1$، أي مجسم زائدي ذو طية واحدة
ودوراني حول المحور $\R(1,1,1)$.

\textbf{18.} بإتمام المربعات: $(x-1)^2 + (y+2)^2 - (z-1)^2 +
(-1 - 4 + 1 + 4) = 0$، أي
\[
(x-1)^2 + (y+2)^2 = (z-1)^2 :
\]
وقد انعدم الثابت --- أي \emph{مخروط} دائري قائم رأسه (ومركزه الوحيد)
$(1, -2, 1)$ ومحوره موازٍ للمحور $Oz$.

\textbf{19.} للجزء التربيعي $x^2 + 4xy + y^2$ المصفوفة $\begin{psmallmatrix}1 & 2\\ 2 & 1\end{psmallmatrix}$ (في المستوي
$xy$)، وقيمتاها الذاتيتان $3$ (على $(1,1)$) و$-1$ (على
$(1,-1)$): فمع $u = \frac{x+y}{\sqrt2}$، $v =
\frac{x-y}{\sqrt2}$ يساوي $3u^2 - v^2$، ويكون \omterm{pb:b2:surfaces:1}{السطح من الدرجة
الثانية}
\[
z = \tfrac32u^2 - \tfrac12v^2 :
\]
أي مجسم مكافئ زائدي (فرتبة $A$ هي $2$ وللمتجهة $b$ مركبة
على $\ker A = \R e_z$: فلا مركز).

\textbf{20.} تحوّل الحركة الصلبة معطيات المعادلة وفق $A \mapsto P^{\mathsf T}\!AP$ (\omterm{def:b2:reduction:eigen}{بالقيم
الذاتية} نفسها)، وليس للصيغ الموحدة أي حرية متبقية عدا تبديل
الإحداثيات وضرب المعادلة كلها في عدد ($> 0$ للحفاظ على الكتابة):
ومنه يتطابق شكلان موحدان مركزيان إلى غاية تقايس إذا وفقط إذا تطابقت
قائمتا المعاملات إلى غاية تبديل وعامل موجب مشترك --- ومن أجل المجسم
الإهليلجي، إذا وفقط إذا تطابقت أنصاف المحاور $(a, b, c)$. أما \omterm{def:b2:affine:subspace}{أفينيا}،
فيمكن أيضا ضرب كل إحداثية على حدة ($z_i \mapsto z_i\sqrt{\abs
{\lambda_i}}$)، وهو ما يمحو \omterm{def:b2:reduction:eigen}{القيم
الذاتية} ولا يبقي إلا إشاراتها: فيصير كل مجسم إهليلجي $u^2 + v^2 + w^2 = 1$، أي
الكرة. وتضمن مبرهنة سيلفستر في القصور الذاتي (\cref{thm:b2:quadratic:sylvester}) بقاء البصمة
عند أي تغيير خطي قابل للقلب: فالأنماط الأفينية في السؤال 15 متمايزة
حقا.

\textbf{21.} نوسّم المستوي \omterm{def:b2:affine:subspace}{أفينيا}: $M = P + su +
tv$. عندئذ يكون $q(P + su + tv)$ كثير
حدود درجته $\leq 2$ بدلالة $(s, t)$ (بنشر الصورة التربيعية بثنائية
الخطية)، ومنه فإن المقطع $\{q = 0\}$ مقطع مخروطي من المستوي، وقد يكون
منحلا. ومن أجل $z = xy$ والمستوي $z = c$: $xy = c$، وهو قطع زائد من أجل
$c \neq 0$، ومن أجل $c = 0$ هما المستقيمان الإحداثيان --- أي زوج
المولّدين المارّين بالمبدأ.

\textbf{22.} \emph{المجسم الإهليلجي}: محدود، فلا يحوي أي مستقيم.
\emph{المجسم الزائدي ذو الطيتين} $\frac{z^2}{c^2} -
\frac{x^2}{a^2} - \frac{y^2}{b^2} = 1$: نقصر على $p + tv$؛ فيجب أن
ينعدم معامل $t^2$ وهو $\frac{v_3^2}{c^2} - \frac{v_1^2}{a^2}
- \frac{v_2^2}{b^2}$، ومنه $v_3 \neq 0$ (وإلا $v =
0$)؛ ويعطي
معامل $t$ المقدار $\frac{p_3v_3}{c^2} =
\frac{p_1v_1}{a^2} + \frac{p_2v_2}{b^2}$، وتعطي متراجحة كوشي--شوارتز
\[
\frac{p_3^2}{c^2}
= \frac{c^2}{v_3^2}\Bigl(\frac{p_1v_1}{a^2} +
\frac{p_2v_2}{b^2}\Bigr)^2
\leq \frac{c^2}{v_3^2}\Bigl(\frac{p_1^2}{a^2} +
\frac{p_2^2}{b^2}\Bigr)\frac{v_3^2}{c^2}
= \frac{p_1^2}{a^2} + \frac{p_2^2}{b^2},
\]
ومنه فإن الحد الثابت $\leq 0 \neq 1$: فلا مستقيم.
\emph{المجسم المكافئ الإهليلجي} $z = \frac{x^2}{a^2} +
\frac{y^2}{b^2}$: يفرض معامل $t^2$ أن $v_1 = v_2 =
0$،
وعندئذ تكون المعادلة خطية غير ثابتة بدلالة $t$: فلا مستقيم.
\emph{المخروط} $x^2 + y^2 = z^2$: يحقق مستقيم واقع عليه $q(p) =
q(v) = B(p, v) = 0$ من أجل صورة
لورنتز؛ والتساوي في متراجحة كوشي--شوارتز المستوية $\abs{p_1v_1 + p_2v_2} = \abs{p_3v_3} =
\sqrt{(p_1^2 + p_2^2)(v_1^2 + v_2^2)}$ يفرض $(p_1, p_2)
\parallel (v_1, v_2)$
ثم $p = kv$: ومنه فإن جميع المستقيمات تمر بالرأس --- أي عائلة واحدة.
\emph{الأسطوانات}: في حالتي الأسطوانة الإهليلجية والأسطوانة المكافئة
يفرض معامل $t^2$ أن $v_1 = v_2 = 0$ (فالمولّدات فحسب)؛ وفي حالة الأسطوانة
الزائدية $\frac{x^2}{a^2} - \frac{y^2}{b^2} = 1$، يؤدي $\frac{v_1}a = \pm
\frac{v_2}b$ مع $v_2 \neq 0$ عبر معامل $t$ إلى
$\frac{p_1^2}{a^2} = \frac{p_2^2}{b^2}$، وهو ما يناقض الحد الثابت $1$: ومنه فالمولّدات الشاقولية
وحدها من جديد. ومنه فإن السطوح من الدرجة الثانية مزدوجة التسطير هي
بالضبط المجسم الزائدي ذو الطية الواحدة والمجسم المكافئ الزائدي.

\textbf{23.} غاوس: $x^2 - 2xy - 2zx = (x - y - z)^2 - y^2 -
z^2 - 2yz$، ومنه
\[
q = (x - y - z)^2 - 4yz = (x - y - z)^2 + (y - z)^2 - (y +
z)^2 :
\]
أي ثلاثة مربعات مستقلة بالإشارات $(+, +, -)$ --- فالبصمة $(2, 1)$، وهي
توافق السؤال 17، دون أي حساب للقيم الذاتية. ويفقد غاوس المعطيات
المترية: فالإحداثيات الجديدة ليست متعامدة، ومنه تضيع \omterm{def:b2:reduction:eigen}{القيم الذاتية}
(أي شكل المجسم الزائدي ومحاوره وأطوالها)؛ ولا يبقى إلا النمط
الأفيني.

\textbf{24.} لتكن $p$ و$n$ و$z$ أعداد \omterm{def:b2:reduction:eigen}{القيم الذاتية}
الموجبة والسالبة والمعدومة، $p + n + z = 3$. وتحصر قاعدة ديكارت المقدار
$p$ بعدد تغييرات الإشارة $V$ في $\chi_A$، وتحصر $n$ بعدد
تغييرات الإشارة $V'$ في $\chi_A(-\lambda)$؛ وعلاوة على ذلك، ينتج كل زوج من
المعاملات المتتالية غير المعدومة تغييرا في أحد كثيري الحدود بالضبط،
ومنه $V + V' \leq 3 - z$ (فالجذور المعدومة مرئية على صورة معاملات أخيرة معدومة).
ومنه $p + n = 3 - z \geq
V + V' \geq p + n$: أي التساوي، ومنه $p = V$ بالضبط --- فيمكن قراءة إشارات
\omterm{def:b2:reduction:eigen}{القيم الذاتية}. ومن أجل $\chi_A(\lambda)
= \lambda^3 - 3\lambda^2 - 9\lambda - 5$: تعطي الإشارات $+,-,-,-$ المقدار $V = 1$،
وللمقدار $\chi_A(-\lambda) = -\lambda^3 - 3\lambda^2
+ 9\lambda - 5$ الإشارات $-,-,+,-$: $V' = 2$. والبصمة $(1,
2)$ --- وهو
منسجم مع التفكيك المضبوط $\chi_A =
(\lambda - 5)(\lambda + 1)^2$ في السؤال 6.

\textbf{25.} الخوارزمية: قطّر الجزء التربيعي تقطيرا متعامدا
(بالمبرهنة الطيفية: وهي الخطوة \omterm{def:b2:powerseries:analytic}{التحليلية} العميقة الوحيدة، وهي التي
تجعل التصنيف \emph{إقليديا})؛ ثم حل $A\Omega = -b$ لتقرر بين الأنماط
المركزية والأنماط المكافئة ولتزيح الجزء الخطي حيث أمكن (وهي هندسة
أفينية)؛ ثم اقرأ النمط من الرتبة والبصمة والثابت (وتضمن مبرهنة
سيلفستر أن هذه ثوابت)؛ وعندما لا يهم إلا النمط الأفيني، يحل \omterm{thm:b2:quadratic:gauss}{اختزال
غاوس} محل المبرهنة الطيفية على حساب المعلومات المترية. وفي $\R^2$
تصنف الآلة نفسها المقاطع المخروطية: القطع الناقص والقطع الزائد
والقطع المكافئ، زائد أزواج المستقيمات، ومستقيم، ونقطة، ومجموعات
خالية. وفي $\R^n$ لا يتغير شيء عدا مسك الحسابات: فالأنماط تُفهرس
ببصمة $A$، وموضع $b$ بالنسبة إلى $\operatorname{im}A$، وثابت واحد --- مع
المصفوفة المحفوفة ذات البعد $(n{+}2)$ في $(A, b, c)$ التي توفر ثابتا
مضغوطا.
\end{solution}
